% !TeX root = ../main.tex
% Original conceptual sequence: two distinct unconnected objects temporarily
% co-move, then move independently; A rotates and exposes an EXISTING face.
% A/B labels denote conceptual identities, not an algorithm's inferred labels.
\begingroup
% Shared local styles for the four original problem illustrations in Section 2.2.
% Loaded inside each figure's group; all meanings remain legible in grayscale.
\linespread{1}\selectfont
\tikzset{
  rpBase/.style={x=1cm,y=1cm,text=black,line cap=round,line join=round,
    font=\songti\fontsize{9.8pt}{12pt}\selectfont,
    every node/.style={align=center,inner sep=0pt,outer sep=0pt}},
  rpTitle/.style={font=\heiti\fontsize{10.3pt}{13pt}\selectfont},
  rpNote/.style={font=\songti\fontsize{9.3pt}{12pt}\selectfont},
  rpFrame/.style={draw=black,line width=0.65pt},
  rpArrow/.style={draw=black,line width=0.8pt,
    -{Latex[length=1.7mm,width=1.15mm]}},
  rpMotion/.style={draw=black,line width=0.95pt,
    -{Latex[length=1.6mm,width=1.1mm]}},
  rpAction/.style={draw=black,line width=1.5pt,
    -{Latex[length=2.2mm,width=1.7mm]}},
  rpHistory/.style={draw=black!65,line width=0.65pt,
    dash pattern=on 1.4pt off 1.6pt},
  rpFuture/.style={draw=black,line width=0.75pt,
    dash pattern=on 3pt off 2pt},
  rpBoundary/.style={draw=black,line width=1.25pt}
}
\newcommand{\rpPanel}[5]{%
  \draw[rpFrame] (#1,#2) rectangle (#3,#4);
  \path[fill=black!5] (#1+0.01,#4-0.55) rectangle (#3-0.01,#4-0.01);
  \draw[rpFrame] (#1,#4-0.55) -- (#3,#4-0.55);
  \node[rpTitle] at ({(#1+#3)/2},{#4-0.275}) {#5};
}
\newcommand{\rpFooter}[1]{%
  \draw[rpFrame] (0.10,0.10) rectangle (15.90,0.77);
  \node at (8,0.435) {#1};
}
% Hatching is explicitly drawn and clipped, so no additional TikZ libraries
% or color meanings are required. Argument: a closed local-coordinate path.
\newcommand{\rpHatchFace}[1]{%
  \path[fill=black!5] #1;
  \begin{scope}
    \clip #1;
    \foreach \rhx in {-6,-5.8,...,8}{%
      \draw[black,line width=0.35pt] (\rhx,-1) -- ++(3,6);
    }
  \end{scope}
  \draw[rpBoundary] #1;
}

\begin{tikzpicture}[rpBase]
  \path[use as bounding box] (0,0) rectangle (16,5.35);
  \rpPanel{0.10}{1.05}{5.00}{5.20}{短时共同运动}
  \rpPanel{5.55}{1.05}{10.45}{5.20}{后续相对运动}
  \rpPanel{11.00}{1.05}{15.90}{5.20}{已有表面首次显露}

  % Same front-facing rectangular objects; a visible gap means no material bond.
  \node at (2.55,4.24) {$t_1$};
  \draw[rpFrame,fill=black!5] (0.70,2.90) rectangle (1.85,3.75);
  \draw[rpFrame,fill=white] (2.80,2.90) rectangle (3.95,3.75);
  \node at (1.275,3.325) {A};
  \node at (3.375,3.325) {B};
  \draw[rpMotion] (0.85,2.48) -- (1.75,2.48);
  \draw[rpMotion] (2.95,2.48) -- (3.85,2.48);
  \draw[rpFrame] (0.40,1.30) rectangle (4.70,1.85);
  \node[rpNote] at (2.55,1.575) {相似运动不等于同一对象};

  \draw[rpArrow] (5.07,3.05) -- (5.48,3.05);
  \node at (8.00,4.24) {$t_2$};
  \draw[rpFrame,fill=black!5] (6.25,3.10) rectangle (7.40,3.95);
  \draw[rpFrame,fill=white] (8.50,2.35) rectangle (9.65,3.20);
  \node at (6.825,3.525) {A};
  \node at (9.075,2.775) {B};
  \draw[rpMotion] (7.60,3.45) -- (8.20,3.95);
  \draw[rpMotion] (9.80,2.60) -- (10.20,2.10);
  \draw[rpFrame] (5.85,1.30) rectangle (10.15,1.85);
  \node[rpNote] at (8.00,1.575) {长期运动提供归属证据};

  \draw[rpArrow] (10.52,3.05) -- (10.93,3.05);
  \node at (12.20,4.24) {$t_3$};
  % A turns: the newly visible right face belongs to the same existing object.
  \path[draw=black,fill=white,line width=0.65pt]
    (11.45,3.50) -- (12.00,3.90) -- (13.40,3.90) -- (12.85,3.50) -- cycle;
  \path[draw=black,fill=black!5,line width=0.65pt]
    (11.45,2.60) rectangle (12.85,3.50);
  \begin{scope}[shift={(11.45,2.60)}]
    \rpHatchFace{(1.40,0) -- (1.95,0.40) -- (1.95,1.30) -- (1.40,0.90) -- cycle}
  \end{scope}
  \node at (12.15,3.05) {A};
  \draw[rpFrame] (13.65,3.85) rectangle (15.60,4.38);
  \node[rpNote] at (14.625,4.115) {首次可见};
  \draw[draw=black,line width=0.55pt] (13.65,3.95) -- (13.15,3.33);
  \fill[black] (13.15,3.33) circle[radius=0.035];
  \draw[rpFrame] (11.30,1.30) rectangle (15.60,1.85);
  \node[rpNote] at (13.45,1.575) {新观测补入原对象 A};

  \rpFooter{研究关注：连续运动如何支持稳定归属和局部信息累积？}
\end{tikzpicture}
\endgroup
